% =====================================================================
%  Methods — model, inference and objective
%  Updated against sherlock/tools/_models.py (VAE.model, VAE.guide,
%  HardConcreteGate, QRCov) and sherlock/tools/_datasets.py
%  (PerturbMatchingDataset.get_collate_fn).
%
%  CHANGES FROM THE PREVIOUS VERSION (all others left as they were):
%   1. Latent transition is additive (z = z0 + Delta). The previous
%      product-of-experts fusion, Eq. (eq:poe) and the diagonal scale W_s
%      are removed; the transition now has a fixed isotropic variance.
%   2. The background z0 of a perturbed cell is a MONTE CARLO draw: each
%      perturbed cell is paired with a real NTC cell sampled with
%      replacement from its condition's NTC pool, and z0 is the encoding of
%      that cell. Previously z0 was encoded from a single background per
%      cell without the resampling/dedup structure.
%   3. The z0 prior is standard normal in all cases. The condition-dependent
%      prior with learned embeddings (mu^(c), sigma^(c)) is no longer in the
%      model; condition dependence enters through which NTC pool the paired
%      background is drawn from.
%   4. Added the gate row-repulsion regulariser, and the warm-up weights on
%      the KL and gate-sparsity terms.
%   5. Added the reconstruction weights lambda_ntc, lambda_p.
%
%  \eqref{eq:full-elbo} and the equation numbering are preserved so the
%  combinatorial Methods subsection continues to reference them.
% =====================================================================

\section{Methods}

\subsection{Notations}
Let $G$ be the number of genes, $d$ the latent dimension, $P$ the number of perturbation labels, and
$C$ the number of experimental conditions. For a cell $n=1,\dots,N$ we observe: (i) a perturbation
label $p_n \in \{1,\dots,P\}$; (ii) optional condition label $c_n \in \{1,\dots,C\}$; (iii) a count
vector $x^{(\mathrm{p})}_n \in \mathbb{N}^G$ under perturbation, together with a count vector
$x^{(\mathrm{ntc})}_n \in \mathbb{N}^G$ from a non-target control (NTC) cell drawn uniformly at
random, with replacement, from the NTC pool of cell $n$'s own condition. The NTC cell paired with
$x^{(\mathrm{p})}_n$ therefore varies between epochs and serves as a Monte Carlo sample of the
background population rather than as a fixed partner. Let $\theta \in \mathbb{R}^G_{>0}$ denote
Negative Binomial (NB) total-count parameters, one per gene.

\subsection{Generative model}
VAE latent spaces lack identifiability without explicitly regularizing the generative process. In
order to learn interpretable representations, we inject latent geometry structure via a low rank
perturbation covariance and a sparsity gating mechanism.

\subsubsection{Low-rank perturbation covariance and gated shifts}
We introduce a latent matrix $A \in \mathbb{R}^{P \times d}$ that encodes perturbation-specific
shifts in latent space. The rows of $A$ are Gaussian with shared row-covariance $\Sigma_P$:
\begin{align}
  A &\sim \Normal\!\big(0,\ \Sigma_P \otimes I_d\big),
  &
  \Sigma_P \;=\; L L^\top + \sigma^2 I_P,
  \label{eq:row-cov}
\end{align}
where $L \in \mathbb{R}^{P \times r}$ comes from an orthonormal--scaled QR parametrization
$L = Q \diag(s)$ with $Q^\top Q=I_r$ and $s \in \mathbb{R}^r_{\ge 0}$, and $\sigma>0$ is a learned
scalar. The factorization is canonicalized at every evaluation by fixing the sign of each column of
$Q$ and ordering the columns by decreasing $s$, which removes the sign and permutation ambiguity of
the QR decomposition. More specifically,
\[
A \;=\; \mathrm{chol}(\Sigma_P)\,\rho,\qquad
\rho \stackrel{\text{i.i.d.}}{\sim} \Normal(0, I_d)\ \ \text{row-wise}.
\]

To sparsify and select which latent coordinates respond to each perturbation, we use a HardConcrete
gate $W \in [0,1]^{P\times d}$. For perturbation $p$ and latent coordinate $j$,
\begin{align}
  s_{pj} &\sim \text{Concrete}\big(\alpha_{pj}, \tau\big),\qquad
  \bar{s}_{pj} \;=\; s_{pj}(\zeta-\gamma) + \gamma, \\
  W_{pj} &= \mathrm{clip}(\bar{s}_{pj}, 0, 1)\; +\; \big(\bar{s}_{pj} - \mathrm{clip}(\bar{s}_{pj}, 0,1)\big)_{\text{stop-grad}},
\end{align}
with $\tau>0$ temperature and stretch parameters $\gamma<0<1<\zeta$; at evaluation the stochastic
Concrete draw is replaced by its deterministic counterpart $s_{pj} = \sigma(\alpha_{pj}/\tau)$.
The expected L0 usage per gate used in regularization is
\begin{align}
  \pi_{pj}\;\coloneqq\; \mathbb{E}[\1\{W_{pj}>0\}]
  \;=\; \sigma\!\left(\alpha_{pj} - \tau \log\!\frac{-\gamma}{\zeta}\right),
  \label{eq:expected-l0}
\end{align}
where $\sigma(\cdot)$ is the logistic sigmoid and $\alpha_{pj}$ is the learnable logit.

\subsubsection{Background state and latent transition}
Each cell has two latent Gaussian variables:
\begin{itemize}
\item A background latent $z^0_n \in \mathbb{R}^d$, modeling the unperturbed cellular state, with an
isotropic prior
\begin{align}
  z^0_n \,\sim\, \Normal(0, I_d).
\end{align}
Condition dependence is not carried by this prior. It enters instead through the pairing described
in the Notations: the NTC cell whose encoding supplies $z^0_n$ is drawn from cell $n$'s own
condition, so the background is a sample from the appropriate control population by construction.
\item A perturbed latent $z_n \in \mathbb{R}^d$ obtained by displacing the background with the gated
shift of the assigned perturbation. Define
\[
\Delta_{n} \;=\; W_{p_n} \odot A_{p_n}\ \in \mathbb{R}^d,
\]
and take the transition to be additive with fixed isotropic variance $\sigma_z^2$:
\begin{align}
  z_n \;\sim\; \Normal\!\big(z^0_n + \Delta_n,\ \sigma_z^2 I_d\big).
  \label{eq:linear-shift}
\end{align}
\end{itemize}
The gate enters only through $\Delta_n$: it selects and scales the coordinates a perturbation moves
and leaves $z^0_n$ itself unchanged, so background variation is preserved along the coordinates a
perturbation does not act on. NTC cells receive $\Delta_n = 0$ and hence $z_n \approx z^0_n$.

\subsubsection{Monte Carlo background sampling}
Because $z^0_n$ for a perturbed cell is the encoding of a randomly drawn NTC cell rather than of
$x^{(\mathrm{p})}_n$ itself, each minibatch supplies a fresh background draw for every perturbed
cell, so that over training the expectation in the objective is taken over the control population as
a whole rather than over any one background state. Sampling with replacement means one NTC cell may be
paired with several perturbed cells within a batch. The duplicated rows are exact copies, so before
evaluating the background reconstruction term we reduce the batch's NTC rows to the distinct cells
actually drawn; this keeps the NTC likelihood from being counted once per pairing and keeps the cost
of the background decode independent of the perturbed batch size. The same reduction is applied in
the model and in the guide so that the two agree on the number and order of background rows.

\subsubsection{NB decoder and auxiliary classifier}
Given a decoder $g_\phi:\mathbb{R}^d \to \mathbb{R}^G$, let
\begin{align}
  \ell_n &= g_\phi(z_n)\in\mathbb{R}^G, &
  \pi_n &= \softmax(\ell_n)\in\Delta^{G-1}, &
  \mu_n &= \big(\textstyle\sum_{g=1}^G x^{(\mathrm{p})}_{ng}\big)\,\pi_n.
\end{align}
With gene-wise NB total-count parameters $\theta \in \mathbb{R}^G_{>0}$,
\begin{align}
  x^{(\mathrm{p})}_n \sim \NB\!\big(\theta,\ \log(\mu_n{+}\varepsilon)-\log\theta\big),
\end{align}
where $\varepsilon>0$ is a small stabilizer.

Analogously, the background counts use $z^0_n$:
\begin{align}
  \tilde{\ell}_n &= g_\phi(z^0_n),\quad
  \tilde{\pi}_n = \softmax(\tilde{\ell}_n),\quad
  \tilde{\mu}_n = \big(\textstyle\sum_{g=1}^G x^{(\mathrm{ntc})}_{ng}\big)\,\tilde{\pi}_n, \\
  x^{(\mathrm{ntc})}_n \sim \NB\!\big(\theta,\ \log(\tilde{\mu}_n{+}\varepsilon)-\log\theta\big).
\end{align}
The two reconstruction terms carry separate weights $\lambda_{\mathrm{ntc}}$ and $\lambda_{p}$,
allowing the background and perturbed likelihoods to be balanced independently of their differing
row counts.

An auxiliary classifier $h_\psi:\mathbb{R}^d\to\mathbb{R}^P$ predicts the perturbation label from
$z_n$:
\begin{align}
  y_n \mid z_n \;\sim\; \mathrm{Cat}\!\left(\softmax(h_\psi(z_n))\right),
\end{align}
and the model includes a cross-entropy term encouraging $y_n=p_n$.

\paragraph{Priors and stabilizers.}
\begin{itemize}
\item $\rho$ prior: for each row $p$, $\rho_p \sim \Normal(0,I_d)$ and $A$ is formed via
\eqref{eq:row-cov}.
\item NB totals: $\theta=\softplus(\vartheta)+10^{-3}$ with free $\vartheta\in\mathbb{R}^G$.
\item Ridge/stabilizers implemented as \emph{log-density factors}:
  \begin{align}
    \log p(\sigma) &\propto -\tfrac12 (\sigma/0.05)^2, \\
    \log p(L) &\propto -\lambda_{\mathrm{cov}}\|L\|_F^2,
  \end{align}
  with $\lambda_{\mathrm{cov}}>0$.
\end{itemize}

\subsection{Variational Inference}
We use amortized mean-field variational families:
\begin{align}
  q(\rho) &= \prod_{p=1}^P \Normal\!\big(\rho_p \mid m_p,\ I_d\big),
  && m = f_\eta(\text{emb}_p)\in\mathbb{R}^{P\times d},\\
  q(z^0 \mid x^{(\mathrm{ntc})}) &= \prod_{n=1}^N \Normal\!\big(z^0_n \mid \mu^0_n,\ \diag(\sigma^{0\,2}_n)\big),
  && [\mu^0_n,\log\sigma^{0\,2}_n] = e_{\phi_0}(x^{(\mathrm{ntc})}_n),\\
  q(z \mid x^{(\mathrm{p})}) &= \prod_{n=1}^N \Normal\!\big(z_n \mid \mu_n^q,\ \diag(\sigma^{q\,2}_n)\big),
  && [\mu^q_n,\log\sigma^{q\,2}_n] = e_{\phi}(x^{(\mathrm{p})}_n).
\end{align}
Here $f_\eta$ maps the learned perturbation embeddings to row means for $\rho$, while
$e_{\phi_0},e_\phi$ are separate encoder networks for $z^0$ and $z$; keeping them separate means the
perturbed-cell encoder receives reconstruction and classification gradients without the background
encoder's KL term acting on the same weights. For a perturbed cell, $e_{\phi_0}$ is applied to the
paired NTC cell rather than to $x^{(\mathrm{p})}_n$, so the posterior background is an encoding of
real control data; for an NTC cell, $z^0_n$ and $z_n$ coincide and both are given by $e_{\phi_0}$.
Gates $W$ are used \emph{inside the model}; their gradients flow through the Hard Concrete
reparameterization. The expected L0 in \eqref{eq:expected-l0} provides a differentiable surrogate for
sparsity.

\subsection{ELBO and regularization}
Let $x=(x^{(\mathrm{p})}_{1:N},x^{(\mathrm{ntc})}_{1:N})$ and $y=p_{1:N}$.
The evidence lower bound (ELBO) augmented with stabilizers and gating regularizers is
\begin{align}
  \mathcal{L}
  &= \E_{q}\Big[
        \underbrace{\lambda_{p}\log p_\phi\big(x^{(\mathrm{p})}\mid z,\theta\big)}_{\text{NB recon (perturbed)}}
      + \underbrace{\lambda_{\mathrm{ntc}}\log p_\phi\big(x^{(\mathrm{ntc})}\mid z^0,\theta\big)}_{\text{NB recon (background)}}
      + \underbrace{\log p_\psi\big(y\mid z\big)}_{\text{aux CE}}
\\[-2pt]&\qquad\qquad
      + \underbrace{\log p(A\mid L,\sigma)}_{\text{low-rank row-Gaussian}}
      + \underbrace{\log p(\sigma)}_{\text{half-normal energy}}
      + \underbrace{\log p(L)}_{\text{cov ridge}}
      + \underbrace{\log p(\theta)}_{\text{(implicit via parameterization)}}
\\[-2pt]&\qquad\qquad
      + \underbrace{\log p(z \mid z^0, W, A)}_{\text{additive transition, Eq.~\eqref{eq:linear-shift}}}
      + \underbrace{\log p(z^0)}_{\text{standard Gaussian}}
    \Big]
\\&\quad
  - \beta\Big[
    \underbrace{\KL\!\big(q(\rho)\,\|\,p(\rho)\big)}_{\text{row means}}
  + \underbrace{\KL\!\big(q(z^0\!\mid x^{(\mathrm{ntc})})\,\|\,p(z^0)\big)}_{\text{background}}
  + \underbrace{\KL\!\big(q(z\!\mid x^{(\mathrm{p})})\,\|\,p(z\mid z^0,W,A)\big)}_{\text{perturbed}}
  \Big]
\\&\quad
  + \underbrace{\mathcal{R}_{\text{gates}}(W;A)}_{\text{structured sparsity/usage}}
  \ +\ \text{const}.
  \label{eq:full-elbo}
\end{align}
The KL terms are scaled by an annealing weight $\beta$ that is ramped from $0$ to $1$ over the first
epochs of training.

\paragraph{Gating regularizers.}
The model adds the following \emph{log-factor} terms, summed over cells or perturbations as in the
implementation:
\begin{align}
  \mathcal{R}_{\text{gates}}
  &= \kappa\Big(-\lambda_{0}\,\sum_{p,j}\pi_{pj}
   \;-\; \lambda_{1}\,\big\|\sigma(\alpha)\big\|_{1}
   \;-\; \lambda_{2}\,\sum_{p,j} \pi_{pj}\,A_{pj}^2\Big)
\\&\quad
   +\ \lambda_{H}\,\underbrace{\Big(-\textstyle\sum_{j=1}^d \bar{p}_j \log(\bar{p}_j+\varepsilon)\Big)}_{\text{column entropy}}
   \;-\; \lambda_{R}\,\underbrace{\frac{1}{P(P-1)}\sum_{p \neq p'} \Big(\tfrac{W_p^\top W_{p'}}{\|W_p\|\,\|W_{p'}\|}\Big)^{2}}_{\text{row repulsion}},
  \qquad
  \bar{p}_j \;=\; \frac{\sum_{p=1}^P \pi_{pj}}{\sum_{p=1}^P\sum_{j'=1}^d \pi_{pj'}}.
\end{align}
Here $\pi_{pj}$ is \eqref{eq:expected-l0}. The first term is an L0 surrogate encouraging sparsity; the
second is an L1 on gate probabilities; the third discourages large gated shifts in proportion to
expected usage; the fourth promotes balanced usage across latent coordinates via entropy of the
column-wise normalized usage; and the last penalizes the squared cosine similarity between the gate
rows of distinct perturbations, encouraging different perturbations to occupy different latent
coordinates. The sparsity terms carry a weight $\kappa$ that is held at zero until KL annealing has
finished and is then ramped to one, so that gates are not driven to collapse before the latent space
has formed. Coefficients $\lambda_0,\lambda_1,\lambda_2,\lambda_H,\lambda_R$ are tunable
hyperparameters.

\paragraph{Auxiliary cross-entropy.}
With logits $h_\psi(z_n)$ and labels $p_n$, the contribution is
\[
\log p_\psi(y\mid z) \equiv -\lambda_{\mathrm{CE}}\sum_{n=1}^N \mathrm{CE}\!\left(\softmax(h_\psi(z_n)),\,p_n\right),
\]
which appears as a factor in the expectation in \eqref{eq:full-elbo}.

\paragraph{Closed-form terms}
The $\KL$ terms involving Gaussians have closed forms:
\begin{align}
\KL\big(\Normal(\mu^q,\diag(\sigma^{q\,2}))\,\|\,\Normal(\mu^p,\diag(\sigma^{p\,2}))\big)
=\tfrac12\sum_{j=1}^d \Big[
\log \frac{\sigma^{p\,2}_j}{\sigma^{q\,2}_j}
+\frac{\sigma^{q\,2}_j + (\mu^q_j-\mu^p_j)^2}{\sigma^{p\,2}_j}-1\Big].
\end{align}
For $q(\rho)$ vs.\ $p(\rho)=\Normal(0,I)$, $\mu^p=0$, $\sigma^{p\,2}_j=1$.
For $q(z^0)$ vs.\ $p(z^0)=\Normal(0,I)$, likewise.
For $q(z)$ vs.\ the transition prior, $\mu^p = z^0_n + \Delta_n$ and
$\sigma^{p\,2}_j = \sigma_z^2$ from \eqref{eq:linear-shift}.

\paragraph{NB likelihood.}
With gene-wise total-counts $\theta_g$, the decoder uses logits
$\log(\mu_{ng}{+}\varepsilon)-\log \theta_g$; this yields the standard NB log-probability for integer
counts. Gradients flow through $\mu_n$ (via $g_\phi$) and through $\theta$.

\paragraph{Optional components.}
Two further terms are implemented and disabled in all experiments reported here. A gene-space
alignment loss supervises the counterfactual path directly, applying the learned shift to the encoded
NTC backgrounds and penalizing the squared deviation of the decoded effect from the observed
differential-expression profile. A Gene Ontology prior replaces the zero mean of $p(\rho)$ with a
learned embedding of each perturbation's annotated terms, turning the row-wise KL into shrinkage
toward an annotation-informed mean rather than toward the origin.
